\documentclass[11pt,a4paper]{article}
\usepackage{amsmath,amssymb,amsthm}
\usepackage[margin=2.5cm]{geometry}
\usepackage{hyperref}

\newtheorem{definition}{Definition}[section]
\newtheorem{theorem}[definition]{Theorem}
\newtheorem{proposition}[definition]{Proposition}
\newtheorem{corollary}[definition]{Corollary}
\newtheorem{remark}[definition]{Remark}
\newtheorem{conjecture}[definition]{Conjecture}

\title{Hyperseed Formalization:\\
  Can All Human Concepts Be Reduced to Semantic Primitives?}
\author{ProtomegaTron (automated formalization)}
\date{2026-09-19 (deepened v2)}

\begin{document}
\maketitle

\begin{abstract}
We formalize Ben Goertzel's ``Can All Human Concepts Be Reduced to
Semantic Primitives?'' (2022-03-14) within the Hyperseed ontology,
incorporating d-calculus extensions from Hyperseed v2. Semantic primitives
are modeled as base fibers forming a fiber basis for the concept space,
embodiment grounding as fiber anchoring, and the limits of reduction as
gradient singularities and compositional holonomy.
\end{abstract}

\tableofcontents

\section{Primitives as fiber basis}

\begin{definition}[Primitive fiber set --- claims 1, 7]
The \emph{primitive fiber set} $\mathcal{P} = \{p_1, \ldots, p_k\}$
consists of fibers that cannot be decomposed into simpler fibers within
the system. Each $p_i$ is an atomic fiber element.
\end{definition}

\begin{definition}[Fiber composition --- claims 4, 7]
Every compound concept $C$ is a composition of primitives:
\[
  C = f(p_{i_1}, p_{i_2}, \ldots, p_{i_n})
\]
where $f$ is drawn from a set of composition operations
$\mathcal{F} = \{\wedge, \vee, \neg, \forall, \exists, \text{metaphor},
\text{blend}, \ldots\}$.
\end{definition}

\begin{proposition}[Fiber basis completeness --- claims 1, 2, 8]
\label{prop:basis}
The primitives $\mathcal{P}$ form a fiber basis if they span the
concept fiber space $\mathcal{C}$:
\[
  \text{span}_\mathcal{F}(\mathcal{P}) = \mathcal{C}
\]
The basis is \emph{complete} if every concept in $\mathcal{C}$ has a
finite composition in terms of $\mathcal{P}$. If some concepts resist
finite decomposition, the basis is incomplete and must be extended.
\end{proposition}

\section{Embodiment grounding as fiber anchoring}

\begin{definition}[Grounding anchor --- claims 3, 9]
An \emph{embodied primitive} $p_e \in \mathcal{P}$ is anchored to the
physical base space through sensorimotor interaction:
\[
  \text{anchor}: p_e \to B_{\text{sensorimotor}}
\]
Anchored primitives have direct physical referents (MOVE $\to$ motor
action, SEE $\to$ visual perception). They connect the abstract fiber
bundle to the physical base space.
\end{definition}

\begin{definition}[Metaphorical extension --- claim 10]
An \emph{abstract primitive} $p_a$ is constructed from an embodied
primitive $p_e$ via a fiber morphism:
\[
  \mu: p_e \to p_a, \quad \text{where } \mu \text{ preserves structure}
  \text{ while detaching from sensorimotor content}
\]
The metaphorical morphism preserves relational structure (if A relates
to B as $p_e$ does to its referent, then the abstract A relates to
abstract B as $p_a$ does to its referent) but strips the specific
sensorimotor binding.
\end{definition}

\section{d-calculus formulation (Hyperseed v2)}

\begin{proposition}[Primitive basis curvature --- claim 11]
\label{prop:curvature}
The curvature $F_\nabla^\mathcal{P}$ of the primitive fiber basis measures
composition smoothness:
\[
  \|F_\nabla^\mathcal{P}\| \ll 1 \implies \text{compositions well-defined,
    smooth concept formation}
\]
\[
  \|F_\nabla^\mathcal{P}\| \gg 1 \implies \text{compositions ambiguous or
    unexpected (primitive interaction effects)}
\]
High curvature regions identify primitive pairs whose composition is
non-trivial --- e.g., combining GOOD and BAD may have high curvature
(context-dependent, potentially contradictory).
\end{proposition}

\begin{proposition}[Decomposition gradient --- claims 6, 12]
\label{prop:gradient}
The decomposition gradient $\nabla_{\text{decomp}} C$ for concept $C$
points toward its primitive decomposition in fiber space:
\[
  \nabla_{\text{decomp}} C = \sum_i \frac{\partial C}{\partial p_i}\, dp_i
\]
When $\nabla_{\text{decomp}} C$ is well-defined and non-singular, $C$ has
a smooth decomposition. When the gradient is singular:
\[
  \det\left(\frac{\partial^2 C}{\partial p_i \partial p_j}\right) = 0
  \implies C \text{ resists decomposition}
\]
This identifies genuinely primitive or basis-incomplete concepts ---
candidates for extending $\mathcal{P}$.
\end{proposition}

\begin{proposition}[Grounding connection curvature --- claim 13]
\label{prop:grounding}
The grounding connection $\Gamma_{\text{ground}}$ between abstract and
embodied fiber has curvature measuring metaphorical distortion:
\[
  F_{\Gamma_{\text{ground}}} = 0 \implies \text{abstract concept is a
    perfect lift of embodied one}
\]
\[
  \|F_{\Gamma_{\text{ground}}}\| > 0 \implies \text{metaphorical extension
    introduces genuinely new content}
\]
Non-zero grounding curvature means the abstract concept contains
information not present in the embodied base --- it is ``more than a
metaphor.''
\end{proposition}

\begin{proposition}[Compositional holonomy --- claim 14]
\label{prop:holonomy}
Composing primitives around a conceptual loop
$\gamma: C \xrightarrow{d_1} C' \xrightarrow{d_2} C'' \xrightarrow{d_3} C$
using different decompositions $d_i$ may produce non-trivial holonomy:
\[
  \operatorname{Hol}_\gamma(\nabla) \neq \mathrm{id}
  \implies \text{decomposition of } C \text{ is ambiguous}
\]
Zero holonomy = unique canonical decomposition; non-trivial holonomy =
multiple valid decompositions, path-dependent meaning.
\end{proposition}

\begin{proposition}[Cross-cultural parallel transport --- claim 15]
\label{prop:transport}
Parallel-transporting a concept $C$ across cultural fiber bundles
$E_{\text{culture}_1} \to E_{\text{culture}_2}$ using primitives as
connection:
\[
  P_\gamma^{12}(C) = C' \quad \text{where } C' \text{ is } C \text{'s
    ``translation'' in culture 2}
\]
If $\mathcal{P}$ is truly universal, the transport is faithful:
$C' \approx C$. If $\mathcal{P}$ is culture-dependent, the transport
introduces holonomy: $C' \neq C$, and the discrepancy $C' - C$ measures
the cultural specificity of the ``universal'' primitives.
\end{proposition}

\end{document}
